\section{境界の取り扱い（B/T維持）と「B面定義の反射をSへ引き戻す」TMM再帰の説明}

本節では，(1) 「境界（EML/EBL）でスタック反射を定義する」という表現が，ユーザのcavity設計（B/T境界）と混同を生む点を整理し，
(2) B面で定義した反射係数をS参照面へ移す操作は，単なる位相因子ではなくTMM（transfer-matrix method）の再帰式で行う必要があることを，
同じ粒度の数式詳細で説明する．
本節ではユーザの設計方針に合わせて，cavityの幾何ゲージ境界は
\[
B \equiv \text{ITO/pHTL},\qquad T \equiv \text{ETL/cathode1}
\]
を固定し，参照面は発光点の面
\[
S \equiv \text{Source plane}
\]
に固定する．

\subsection{(1) 「境界」は1種類ではない：B/Tはcavity境界，EML/EBLは“評価面”の例}

ユーザのcavity設計での境界はB/Tであり，これは長さ・座標のゲージ
\[
(L_{\mathrm{cav}}, z_{\mathrm{cav}})
\]
を定義するための幾何学的境界である．すなわち，
\[
L_{\mathrm{cav}} = \text{(BからTまでのcavity長)},\qquad
z_{\mathrm{cav}} = \text{(Bからのcavity座標)}
\]
はB/Tの選び方に依存する．このB/Tを変えることは「ゲージを変える」ことに等しく，既存コードの安定性・互換性を損ね得る．

一方，反射係数の議論で出てくる「どこで反射を定義するか」は，cavityゲージ境界とは別の自由度であり，
任意の面（評価面）で定義できる．教科書的DGFではsourceが一様層（例：EML）内にあるため，
source層の上下端（例：EML/EBLやEML/ETL）で上下スタック反射を定義するのが自然，というだけである．
しかしユーザがB/T境界を維持する方針なら，評価面をB/Tに取ることも可能である．
重要なのは最終的に「S参照面での有効反射係数」
\[
r_{u,S}(u,\lambda),\qquad r_{d,S}(u,\lambda)
\]
（上側/下側の有効反射）を正しく得ることであり，それを得るための途中の評価面として
EML/EBLを使うかBを使うかは実装選択である．本節ではユーザの設計通り，評価面としてもB/Tを採用し，
B面で定義した鏡反射を多層を介してSへ引き戻す操作を説明する．

\subsection{前提：角スペクトル（$u$）と各層の$k_z$}

真空波長$\lambda$，真空波数$k_0$を
\[
k_0(\lambda)=\frac{2\pi}{\lambda}
\]
とする（コード整合のため$\lambda$はnm，$k_0$は1/nmとして扱ってよい）．
面内波数$k_\parallel$を無次元化して
\[
u \equiv \frac{k_\parallel}{k_0},\qquad k_\parallel = u k_0
\]
とする．層$j$の複素屈折率を$n_j(\lambda)$とし，その層での$z$方向波数を
\[
k_{z,j}(u,\lambda)=k_0(\lambda)\sqrt{n_j(\lambda)^2-u^2}
\]
と定義する．枝は発散解を避けるため
\[
\Im\!\left[k_{z,j}(u,\lambda)\right]\ge 0
\]
となるように選ぶ．層厚を$d_j$（nm）とすると，伝搬位相は
\[
\delta_j(u,\lambda)=k_{z,j}(u,\lambda)\,d_j
\]
である．

\subsection{(2) B面定義の反射をSへ移す：単なる位相ではなくTMM再帰が必要}

\subsubsection{なぜ位相因子だけでは不十分か}
もしSとBの間が「同一屈折率の一様媒質」1層で距離$d$だけなら，
Bでの反射をSで見た反射へ移す操作は往復位相を掛けるだけで
\[
r_{d,S} = r_B\,e^{2 i k_{z} d}
\]
で済む．これは途中に界面反射がない（連続媒質）ことを意味する．

しかし実際にはS--B間に複数層があり，それぞれの界面でフレネル反射が存在する．
Bで反射して戻ってきた波は途中界面でさらに反射し，再びBで反射し，という多重反射系列を形成するため，
「Bでの反射に位相を掛けるだけ」ではSでの見かけの反射を表せない．
この多重反射系列を正しく畳み込むと，反射は「足し算」ではなく幾何級数の和として閉じ，
最終的に有理式（分数）となる．これがTMM再帰式の本質である．

\subsubsection{最小例：単一薄膜を介して下側反射$r_2$を見る（足し算$\to$幾何級数$\to$分数）}
媒質0（S側）と媒質2（下側スタック側）の間に，厚み$d_1$の薄膜（媒質1）があるとする：
\[
\text{medium 0}\;|\;\text{film 1 (thickness }d_1)\;|\;\text{medium 2 (stack)}
\]
film 1の下側（medium 2側）から見た有効反射を$r_2$とする（これが「B面で定義した鏡反射」を代表する）．
0/1界面の反射係数を$r_{01}$，透過係数を$t_{01},t_{10}$とし，
film内位相を
\[
\delta_1 = k_{z,1} d_1
\]
とする．0側から見える反射は，まず
\[
r_{01}
\]
（界面で1回反射）があり，さらに薄膜を往復して下側で反射して戻る成分
\[
t_{01}t_{10}\,r_2\,e^{2i\delta_1}
\]
が加わる．しかし，この戻り波は0/1界面で再反射し，再び薄膜を往復して下側で反射して戻るため，
無限回の多重反射が起きる．その結果，反射は幾何級数の和となり，
閉形式として
\[
\boxed{
r_0
=
\frac{r_{01}+r_2\,e^{2i\delta_1}}
{1+r_{01}r_2\,e^{2i\delta_1}}
}
\]
となる．ここで$r_0$が0側（S側）から見た「薄膜＋下側スタック全体」の有効反射である．
この式は「位相を掛けるだけ」の形
\[
r_0 \neq r_2 e^{2i\delta_1}
\]
であることを明示している．界面反射$r_{01}$が存在する限り，分母を持つ有理式が不可避となる．

\subsubsection{一般化：多層を1枚ずつ巻き戻すTMM再帰式（TE/TM別）}
薄膜が複数枚ある場合も同様に，最下層側（B面でまとめた鏡反射）から出発して，
層を1枚ずつ上へ「巻き戻す」ことで，任意の上側面での有効反射を得られる．
層$j$とその下の層$j+1$の界面（$j|j+1$）を考え，
層$j+1$以降の全スタックを層$j+1$側から見た有効反射を$r_{j+1}^{(\sigma)}$とする．
このとき層$j$側から見た有効反射$r_j^{(\sigma)}$は
\[
\boxed{
r_j^{(\sigma)}
=
\frac{
r_{j,j+1}^{(\sigma)} + r_{j+1}^{(\sigma)} e^{2i\delta_{j+1}}
}{
1 + r_{j,j+1}^{(\sigma)}\, r_{j+1}^{(\sigma)} e^{2i\delta_{j+1}}
}
}
\]
で与えられる．ここで
\[
\delta_{j+1} = k_{z,j+1} d_{j+1}
\]
であり，偏光$\sigma\in\{s,p\}$（TE/TM）ごとに界面反射$r_{j,j+1}^{(\sigma)}$が異なる．

界面反射はアドミタンス表現で一括に書ける．層$j$の（規格化）アドミタンスを
\[
Y_j^{(s)} \propto k_{z,j},
\qquad
Y_j^{(p)} \propto \frac{n_j^2 k_0^2}{k_{z,j}}
\]
とすると（比例定数は相殺される），
\[
r_{j,j+1}^{(\sigma)}=
\frac{Y_j^{(\sigma)}-Y_{j+1}^{(\sigma)}}{Y_j^{(\sigma)}+Y_{j+1}^{(\sigma)}}.
\]
この再帰をB面の評価で得た反射$r_B^{(\sigma)}(u,\lambda)$から開始し，
Sに到達するまでS--B間に存在する層列（pHTL，Rprime，HTL，EBL，EML下半分，など）を順に巻き戻せば，
S参照面での下側有効反射
\[
r_{d,S}^{(\sigma)}(u,\lambda)
\]
が得られる．同様にT面で定義した反射$r_T^{(\sigma)}(u,\lambda)$を出発点として，
S--T間の層列（ETL，EML上半分，など）を巻き戻すことで，上側有効反射
\[
r_{u,S}^{(\sigma)}(u,\lambda)
\]
が得られる．

\subsubsection{現状proxy（位相のみ）との関係：界面反射をゼロとした極限}
途中の界面反射が全て無視できる（連続媒質の極限）なら
\[
r_{j,j+1}^{(\sigma)}\to 0
\]
となり，再帰式は
\[
r_j^{(\sigma)} \approx r_{j+1}^{(\sigma)} e^{2i\delta_{j+1}}
\]
へ退化する．これは「位相因子を掛けるだけ」という現行proxyの操作に一致する．
したがって，ユーザが求める（2）の意味で「Sから見た層群全部含む有効反射」を実現するには，
位相のみの引き戻しではなく，上のTMM再帰をS--BおよびS--Tの層列に対して適用する必要がある．
